\documentclass[aps,prl,twocolumn,showpacs,superscriptaddress]{revtex4-2}
\usepackage{amsmath,amssymb,booktabs,xcolor,hyperref,amsthm}
\newtheorem{theorem}{Theorem}

\begin{document}

\title{Appendix KA3: The BCT Bond Angle Theorem --- Full Algebraic Derivation\\
Proof of $r_{\rm tet}(1+r_{\rm tet})=1/8$ and the complete rational
series of sp$^3$ bond angles}

\author{Michel Robert Cabri\'e}
\affiliation{Independent Researcher, Victoria, Australia}
\email{ZeroFreeParameters@gmail.com}
\date{20 March 2026}

\begin{abstract}
This appendix provides the complete algebraic derivation of the BCT 
bond angle theorem, supporting Letter~80. The central result is the 
identity $r_{\rm tet}(1+r_{\rm tet}) = 1/8$ (exact), proved by the 
difference of squares $(\sqrt{6}-2)(\sqrt{6}+2) = 2$. This gives 
the universal sp$^3$ formula $\cos\theta = -(1-n_{\rm lp}/8)/3$, 
generating the rational series: CH$_4 = \arccos(-1/3) = 109.471^\circ$ 
(exact); NH$_3 = \arccos(-7/24) = 106.958^\circ$; 
H$_2$O $= \arccos(-1/4) = 104.4775^\circ$ (obs $104.4776^\circ$, 
error $-0.0001\%$).
\end{abstract}

\maketitle

\section{The Key Identity}

\begin{theorem}
$r_{\rm tet}(1+r_{\rm tet}) = 1/8$, exactly.
\end{theorem}

\begin{proof}
With $r_{\rm tet} = (\sqrt{6}-2)/4$:
\begin{align}
1 + r_{\rm tet} &= \frac{4+\sqrt{6}-2}{4} = \frac{\sqrt{6}+2}{4}\\
r_{\rm tet}(1+r_{\rm tet}) &= \frac{(\sqrt{6}-2)(\sqrt{6}+2)}{16}
= \frac{6-4}{16} = \frac{2}{16} = \frac{1}{8}
\end{align}
The difference of squares $(\sqrt{6})^2 - 2^2 = 2$ completes 
the proof. \qed
\end{proof}

\section{Physical Origin}

In BCT (Letter~78), lone pairs are closed Hopf multiplets ($H=0$) 
confined to the central nucleus. Their effective angular reach in 
the molecular frame is:
\begin{equation}
    r_{\rm tet} + r_{\rm tet}^2 = r_{\rm tet}(1+r_{\rm tet}) = \frac{1}{8}
\end{equation}
where $r_{\rm tet}$ is the primary Hopf mode radius and $r_{\rm tet}^2$ 
is the Josephson self-coupling correction. Each lone pair compresses 
$\cos\theta$ by $1/8$.

\section{The Rational Series}

\begin{center}
\begin{tabular}{lccll}
\toprule
$n_{\rm lp}$ & $\cos\theta$ & $\theta$ & Molecule & Error \\
\midrule
0 & $-1/3$ & $109.471^\circ$ & CH$_4$ & Exact \\
1 & $-7/24$ & $106.958^\circ$ & NH$_3$ & $-0.78\%$ \\
2 & $-1/4$ & $\mathbf{104.4775^\circ}$ & \textbf{H$_2$O} & $\mathbf{-0.0001\%}$ \\
3 & $-5/24$ & $102.025^\circ$ & F$_2$O & theoretical \\
\bottomrule
\end{tabular}
\end{center}

\section{Numerical Verification --- H$_2$O}

\begin{align}
r_{\rm tet} &= (\sqrt{6}-2)/4 = 0.11237243569\ldots\\
r_{\rm tet}(1+r_{\rm tet}) &= 0.12500000000\ldots = 1/8\\
\theta_{\rm H_2O} &= \arccos(-1/4) = 104.47751219\ldots^\circ\\
\theta_{\rm obs} &= 104.47760^\circ\quad\text{(NIST)}\\
\text{Error} &= -0.000084\%
\end{align}

The H$_2$O bond angle is $\arccos(-1/4)$ --- an exact rational 
fraction following from the difference of squares identity on 
$r_{\rm tet} = (\sqrt{6}-2)/4$.

\begin{thebibliography}{99}
\bibitem{L78} M.~R.~Cabri\'e, \textit{BCT Letter 78}, Zenodo (2026).
\bibitem{L80} M.~R.~Cabri\'e, \textit{BCT Letter 80}, Zenodo (2026).
\end{thebibliography}

\end{document}
